% Test file to verify theorem numbering fix
% This document should be used with jbook or jreport class which support chapters
\documentclass{jbook}
\usepackage{kuis-cover}

% Define theorem environment numbered by chapter
\newtheorem{theorem}{定理}[chapter]  % 定理 = Theorem in Japanese
\newtheorem{lemma}[theorem]{補題}     % 補題 = Lemma in Japanese (shares counter with theorem)
\newtheorem{definition}{定義}[chapter] % 定義 = Definition in Japanese (has its own counter)
% Note: lemma shares numbering with theorem, while definition has its own counter.
% This demonstrates both numbering styles work correctly with the fix.

\begin{document}

\chapter{第一章}  % Chapter 1

\section{はじめに}  % Introduction

\begin{theorem}
これは定理です。番号は「定理1.1」と表示されるべきです。
% This is a theorem. It should display as "定理1.1" (Theorem 1.1)
\end{theorem}

\begin{theorem}
二番目の定理です。番号は「定理1.2」と表示されるべきです。
% Second theorem. It should display as "定理1.2" (Theorem 1.2)
\end{theorem}

\begin{lemma}
補題です。番号は「補題1.3」と表示されるべきです。
% A lemma. It should display as "補題1.3" (Lemma 1.3)
\end{lemma}

\begin{definition}
定義です。番号は「定義1.1」と表示されるべきです。
% A definition. It should display as "定義1.1" (Definition 1.1)
\end{definition}

\chapter{第二章}  % Chapter 2

\section{別の章}  % Another chapter

\begin{theorem}
第二章の定理です。番号は「定理2.1」と表示されるべきです。
「定理 Chapter 2.1」とは表示されません。
% Theorem in chapter 2. It should display as "定理2.1" (Theorem 2.1)
% NOT as "定理 Chapter 2.1"
\end{theorem}

\begin{definition}
第二章の定義です。番号は「定義2.1」と表示されるべきです。
% Definition in chapter 2. It should display as "定義2.1" (Definition 2.1)
\end{definition}

\end{document}
